\section{Related Work}
\label{sec:related}

HNSW~\cite{malkov2020hnsw} answers a query by greedy traversal of layered proximity graphs with a bounded candidate
list, whose size (\code{ef\_search}) controls the speed--recall trade-off. Li et al.~\cite{li2025attribute} classify
filtered-ANN methods into pre-, post- and joint filtering and show that the best method depends on selectivity.
Filtered-DiskANN~\cite{gollapudi2023filtered} builds label-aware graphs, and ACORN~\cite{patel2024acorn} traverses the
predicate-induced subgraph of a denser HNSW-style graph; both change the index, whereas we choose between the indexes
PostgreSQL already provides. Amanbayev et al.~\cite{amanbayev2026filtered} compare filtering in FAISS, Milvus and
pgvector and observe that pgvector's planner sometimes prefers approximate scans where exact scans would give perfect
recall. The planner bases this choice on row estimates from the histograms in \code{pg\_stats}~\cite{pgdocs-rowest},
which our router reuses. What is missing is an evaluation of how far a \emph{single calibrated threshold} on these
estimates gets relative to fixed strategies and the default planner, under explicit quality constraints and with
held-out testing.
